\textbf{Complete exploratory screen:} 45 fits, nine participants and seed 0.
All percentages below are participant-equal balanced accuracy on training or
validation data. No test-set performance is claimed.
\begin{center}
\begin{tabular}{@{}lrrrr@{}}\toprule
Model & Parameters & Train (\%) & Validation (\%) & Log loss \\ \midrule
EEGNet--Transformer & 11,092 & 74.69 & 63.71 & 1.029 \\
Spectral dense control & 14,580 & 78.73 & 50.97 & 1.435 \\
Spectral Tucker (A) & 10,340 & 65.17 & 45.47 & 1.625 \\
Channel--time tensor (B) & 11,740 & 76.53 & 44.21 & 1.680 \\
Factorized temporal (C) & 11,000 & 55.85 & 40.82 & 1.394 \\
\bottomrule\end{tabular}\end{center}
\textbf{Accuracy verdict.} Keep the EEGNet--Transformer reference for the
full-input classification goal. None of these concrete tensor versions, nor
the spectral dense control, passes the prespecified accuracy screen against it.
The frozen spectral Tucker model loses 18.24 percentage points to the reference.
This rejects these configurations as accuracy replacements under this protocol;
it does not establish that every tensor encoder is ineffective.
\begin{center}\begin{tabular}{@{}lrrr@{}}\toprule
Matched contrast & Gain (pp) & 95\% interval (pp) & Positive \\ \midrule
Factorized C minus B & -3.39 & [-6.63, -0.46] & 3/9 \\
Frozen A minus spectral dense & -5.50 & [-10.49, -0.87] & 2/9 \\
\bottomrule\end{tabular}\end{center}
The restricted temporal factorization saves 740 parameters but reduces accuracy
in the matched comparison. C has lower log loss than B (1.394 versus 1.680),
so the accuracy result should not be mistaken for uniformly worse probability
predictions. Neither raw tensor variant is a strong classifier in this screen.
\subsection*{Availability tradeoff}
\begin{center}
\begin{tabular}{@{}lrrrrr@{}}\toprule
Model & Full & Static 16 & Dynamic 16 & Static 6 & Dynamic 6 \\ \midrule
EEGNet--Transformer & 63.71 & 35.88 & 34.59 & 31.11 & 31.24 \\
Spectral dense control & 50.97 & 37.33 & 38.86 & 29.14 & 28.67 \\
Spectral Tucker (A) & 45.47 & 36.98 & 36.71 & 33.53 & 34.08 \\
Channel--time tensor (B) & 44.21 & 29.85 & 27.78 & 25.33 & 25.86 \\
Factorized temporal (C) & 40.82 & 29.58 & 28.24 & 24.98 & 25.36 \\
\bottomrule\end{tabular}\end{center}
Frozen spectral Tucker is worth retaining as a robustness hypothesis. With
six retained channels it exceeds the spectral dense control by 4.39 pp
(static; interval [0.64, 8.27], 7/9 positive) and 5.40 pp (dynamic;
[1.41, 10.58], 8/9 positive). Relative to EEGNet--Transformer, its corresponding
gains are 2.42 pp ([0.20, 4.72], 7/9) and 2.84 pp ([$-0.08$, 5.77], 6/9).
These secondary condition-specific intervals are unadjusted and exploratory.
The approximately 34\% severe-loss accuracy, single training seed and full-input
penalty do not establish a strong or generally superior robust classifier.
Frozen versus dense also changes supervision and capacity, so the difference
is not uniquely attributable to the tensor decomposition.
\newpage
\subsection*{Recommended next hypothesis, not a measured result}
Keep the local spectral representation and observed-entry ridge inference,
but test classification-driven fine-tuning of its factors through a
differentiable ridge solve. Initialize frozen and trainable variants from the
same fitted factors and use identical ranks, temporal head and training masks.
This would test whether reconstruction-only fitting discards label information
while preserving the original availability mechanism. It is an untested,
post-screen hypothesis; improvement is not guaranteed. Introducing mixed
22/16/6-channel training would require a separate matched training-regime study.
\subsection*{Evidence and checks}
The implementation passed 17 new model/runner tests and all 72 existing
encoder-candidate tests. Tests include hidden NaNs, temporal patch isolation,
training-only frozen calibration, matched initial tensor operators, optimizer
gradient paths, deterministic repeated fits, selection ties and corrupt-artifact
rejection. Synthetic loss reduction is a software check only.
All 45 selected checkpoints passed exact saved-prediction reload checks.
The independent audit passed for all nine participants. It rebuilt masks,
replayed all 945 validation-condition evaluations, and independently recomputed
participant averages and bootstrap contrasts. The saved audit accompanies the
machine-readable report.
Reports: \code{results/tensor-temporal-screen-v1/report/summary.json},
\code{scores.csv}, \code{robustness.csv}, \code{contrasts.csv}, and
\code{independent_audit.json}. The independent audit can be rerun with
\verb|.venv/bin/python scripts/audit_tensor_temporal_screen.py|.
Its output defaults to \code{/tmp/agfl-tensor-temporal-independent-audit.json}.
\par Study identity:\par
\code{6c3d8edc5a13eefbd378282d593165c2d8ec9a6c3f1fd6aa81655d75c7749a8f}.
\par
The cohort used four independent participant processes, each with one PyTorch
thread; per-fit wall time includes contention and is not a controlled speed benchmark.
